% =====================================================================
%  sec08 -- Errata ledger
% =====================================================================
\section{Errata ledger: eight corrections to the corpus}\label{sec:errata}

A verdict paper is only as good as its ledger. The eight entries below are places where the
internal corpus stated something that later work, in the same corpus, showed to be wrong,
mis-quantified, or mis-attributed. Each entry gives the original statement, the correction, and
the file that decides it. No entry is a change of a \LT{} verdict in this paper; two of them
remove the word ``statistical interval'' from numbers that were being quoted as if they were
certificates, and one of them is a missing file that was not missing.

\vspace{0.2em}
{\footnotesize
\begin{longtable}{@{}L{0.4cm}L{3.7cm}L{6.3cm}L{3.9cm}@{}}
\toprule
\textbf{\#} & \textbf{Original statement} & \textbf{Correction} & \textbf{Decided by} \\
\midrule
\endhead
1 & The $L=2$ decompositions were quoted with Gaussian errors: $\Delta=+0.3942\pm0.0258$
    ($6.0\sigma$), $G''=+1.696\pm0.293$ ($5.8\sigma$), $\Delta_{\mathrm{dyn}}=+0.562\pm0.277$
    ($2.0\sigma$)
  & $E[\kappa^{2}]=\infty$: with $\kappa\le2/\delta^{2}$ and a shell volume $\propto\delta^{3}$,
    $E[\kappa^{2}]\le4C\int_0^{\delta^{*}}\mathrm{d}\delta/\delta=\infty$ (log divergence), and
    the measured tail index is $\alpha\approx1.97\le2$, so the sample variance does not converge
    and no $k\sigma$ interval is asymptotically valid. Withdraw the interval; keep only the
    qualitative agreement of the three estimators in sign
  & \LT{} \code{T04-L2ANALYTIC.md} \S6.2; \code{R04-L2CERT.md}:109; \code{O24-L2-PUSH.md}
    (the measured tail) \\
2 & ``$H1$ must hold for all $a$'', with $a$ written for the coarse amplitude and scanned in
    $\inf_a\beta^{*}(a)=1.8729$
  & There are two different symbols. $a$ is the \emph{lattice spacing} (the cutoff) and $a_0$ is
    the \emph{coarse-field amplitude}. $H1$'s quantifiers are $(L,a)$ and $v$ only; the
    conditioning point is pinned at $a_0=0$. So $\inf_{a_0}$ is not in $H1$'s quantifiers, and the
    uniform-in-$a_0$ strengthening was already a known \LN{}: it is a quantitative cross-section
    of an established no-go, not a new counterexample
  & \LT{}/\LN{} \code{O28-INFA.md} \S0.1--\S0.3; \code{ch15-uvface.tex}:46 \\
3 & $R-1\propto L$: the ratio $R$ ``diverges with volume'', read off a table going
    $R_1=0.337\to R_4=-1.472$
  & Inside any single replicated normalisation $S_L=LS_1$ and $D_L=LD_1$, so $R_L\equiv R_1$
    \emph{exactly} --- the ratio is $L$-independent by arithmetic, not approximately. The table's
    drift is an artefact of mixing two conventions (numerator at fixed block amplitude,
    denominator at another). The real defect of $R$ is that it is a ratio and carries no scale
  & \LT{} \code{T10-C5.md} \S0(c)2 \\
4 & The file \code{GKS-GAUGE.md} ``does not exist'' and its result is unavailable
  & It exists (6841 bytes) and its content is used in the corpus: at $\beta=2.5$, $\mathrm{SU}(3)$
    the character coefficients $c_r(\beta)$ alternate in sign, so the positive-character cone
    exists only for $\beta<\beta_1\approx1/48$ and does not extend to the physical point; the
    minimal-frustration plaquette has $\Cov(G,\Phi')=0.22\beta^{2}>0$
  & \LN{}/\LE{} \code{O04-NONAB-GAP.md}:62, :114, :155; file \code{GKS-GAUGE.md} \\
5 & An instruction was issued to translate \code{ch18} from Chinese to English
  & \code{ch18} was English from the start: measured Han characters $=0$ and CJK punctuation
    $=0$; the chapter title is an English sentence and the file is $637$ lines. The premise of the
    instruction was wrong, and the workflow cost is recorded as the error, not the translation
  & \LE{} \code{LOOP-RUNBOOK.md} \S11.49 (``my premise was wrong in one place: \code{ch18} was always
    English'') \\
6 & $C=1.7454$ \LT{} (one source) versus $C=10.05$ (another), with the discrepancy attributed to a
    range/beta-freezing dispute
  & Both are branch artefacts. $1.7454=0.2971\cdot(2.43)^{2}$ is the \emph{range} branch evaluated
    at the crossing point; $10.05$ comes from dropping the $L=4$--$8$ centre chain and
    recomposing with the old four-lemma set (the same machine with $C_1=1$ and $2/B^{2}$ gives
    $6.46$, not $10.05$). The correct value is $C=3.9512=2.25\cdot B^{*2}$ with $B^{*}=1.3252$
  & \LT{}/\LC{} \code{Q16-CONFLICT.md} \S0, \S4.1--\S4.2 \\
7 & The $\beta=6$, $L=2$ result was $+0.5913$ at $1.3\sigma$, and later at $4.1\sigma$
  & The old $\sigma_\tau$ applies $\tau_F$ to the fluctuation of the \emph{sample variance}
    $z=(F-\bar F)^{2}$, whose decorrelation time is $\tau_z\approx\tau_F/3.5$; using $\tau_F$
    overestimates that block by $1.83\times$. The correct object is $\tau_z$, after which the same
    data give $7.3\sigma$; and at $\beta=6$ the variance block is $96.9\%$ of $\sigma^{2}_{\tau}$,
    so the bottleneck is $F$-fluctuation, not the $\kappa$ tail
  & \LT{}/\LE{} \code{Q13-BETA6.md} \S2, \S2.1 (the $\tau_F$ misuse), \S2.3 ($96.9\%$) \\
8 & ``2D Yang--Mills has $G''>0$ for all $t$'', quoted as evidence that the failure of $G''\ge0$
    is a four-dimensional phenomenon
  & Withdraw: in 2D Yang--Mills $G''>0$ on all of $\varphi\in(0,\pi)$ \emph{if and only if}
    $t\ge t_c=1.3439$; below the threshold there is a window where $G''<0$ (at $t=1$:
    $[1.271,1.871]$, minimum $-2.87$). Also, the 2D and $P_4$ cases are vacuous for the acid test
    ($\Var=E[H]=0$, all margin from $\Wg$), so their global convexity carries no interacting
    content
  & \LE{} \code{O02-LOGCONCAVE.md} \S0 (the table and the final two bullets) \\
\bottomrule
\end{longtable}
}

\subsection{What these eight entries do and do not change}

None of the eight changes a \LT{} verdict of this paper. Two of them (entries 1 and 7) change the
\emph{epistemic status} of numbers that were being read as statistical intervals: entry 1 says
that a $6.0\sigma$ and a $5.8\sigma$ are not intervals at all, because the fourth moment of
$\kappa$ diverges; entry 7 says that a $1.3\sigma$ was a formula mismatch and that the same data
support $4.1\sigma$ (old formula) or $7.3\sigma$ (corrected), so the number was never a property
of the data alone. Two of them (entries 2 and 8) are quantifier corrections: a symbol collision
that let a known no-go be re-read as a new counterexample, and a conditional convexity result that
had lost its condition. Two (entries 3 and 6) are bookkeeping of normalisation: a ratio that was
thought to scale and does not, and a constant that was the value of the wrong branch at the
crossing point. One (entry 5) records a workflow error whose only trace is the instruction itself.
One (entry 4) is a missing file that was not missing --- the file is
\code{research/unconditional/GKS-GAUGE.md}, $6841$ bytes --- and the most embarrassing kind,
because the result it contains (the positive-character cone stops at $\beta\approx1/48$) is
load-bearing for \S\ref{sec:dead} items 8--9.

\src{\code{T04-L2ANALYTIC.md} \S6.2 and \code{:180} (the two new registry entries: ray-wise
quadratic homogeneity, and $E[\kappa^{2}]=\infty$); \code{O28-INFA.md} \S0; \code{T10-C5.md}
\S0(c); \code{O04-NONAB-GAP.md}:62, :114; \code{LOOP-RUNBOOK.md} \S11.49; \code{Q16-CONFLICT.md}
\S0; \code{Q13-BETA6.md} \S2; \code{O02-LOGCONCAVE.md} \S0.}

\subsection{Two further corrections that belong to the certificate, not to the corpus}
\label{sec:errata2}

For completeness, and because \S\ref{sec:cert} is the only place in this paper where a number was
withdrawn, the two certificate-level corrections are repeated here in ledger form rather than left
buried in \S\ref{sec:grading}. (i) The first general-$\theta$ implementation printed
$\operatorname{gap}(0)=0.729924127576750$ with an enclosure of width $1.4\times10^{-10}$; the truth
is $0.729924128620548449$, higher by $1.044\times10^{-9}$, and the enclosure does not contain it,
because the quotient lost the tails of $Z,N_A,N_B,N_S$ (which are themselves valid enclosures).
\LT{}/\LE{} (ii) The same defect exceeds the $\theta=\pi/2$ enclosure by $50\times$; the sign
verdict is unaffected in both cases, and the numbers used in \S\ref{sec:cert} are the audit's.
Additionally, the $\mathrm{SU}(3)$ value $\operatorname{gap}=+0.1017$ is \LE{} only, is
weight-dependent (pure Haar gives $-0.2358$), and its $\Var(\Phi')$ violates the exact identity of
\S\ref{sec:lemmas} lemma 4, so it must be recomputed rather than quoted.

\src{\code{R10-AUDITCERT.md} \S0, \S3; \code{R06-GAPMONO.md} \S4 (the independent recomputation and
its ``worth re-checking whether the tail really entered the quotient'' note);
\code{R12-SU3CORR.md} \S0; \code{S04-VARPHI.md} verdict item three.}

\subsection{Closing status}

Everything displayed in this document is either an identity with a named file, a counterexample
with an exact computation, an interval-decided sign, or a measurement with its scope attached. The
count that matters is unchanged and unimproved: the Clay mass gap for $4$D pure $\mathrm{SU}(N)$
Yang--Mills is open, this paper establishes no progress on it, and the object on which progress
would have to be made is $\fone$ together with I2 (\S\ref{sec:target}) --- not the coarse-mode
curvature, whose two obstructions are the subject of \S\ref{sec:obs1} and \S\ref{sec:obs2}. The
one strict certificate (\S\ref{sec:cert}) is real, and it certifies a $0$-dimensional integral in a
window; it is recorded here because a verdict paper must record what it did prove, not only what
it refuted.
